\section{Related Work}
\label{sec:related}

\subsection{Automated Documentation Generation}

Javadoc~\cite{kramer1999javadoc} and Doxygen~\cite{doxygen2023} generate API documentation from source code comments. Sphinx~\cite{sphinx2023} extends this to narrative documentation with reStructuredText. These tools document APIs but do not generate research papers.

Recent work uses LLMs to generate code documentation~\cite{chen2021codex}. GitHub Copilot~\cite{chen2021evaluating} suggests docstrings based on function signatures. These systems focus on code-level documentation rather than paper-level research contributions.

\subsection{Mining Software Repositories (MSR)}

MSR research extracts insights from software development artifacts~\cite{hassan2008road}. Studies analyze commit messages~\cite{mockus2000identifying}, issue trackers~\cite{bird2009promises}, and code review systems~\cite{rigby2013convergent} to understand developer behavior and software evolution.

Our system extends MSR by generating research papers directly from artifacts rather than extracting data for manual analysis. We automate the full pipeline from commit history to publication-ready LaTeX.

\subsection{Automated Scientific Writing}

Recent work explores LLM-based scientific writing assistance. Elicit~\cite{elicit2023} helps researchers find and summarize papers. Semantic Scholar's TLDR~\cite{cachola2020tldr} generates paper summaries. These systems assist human authors rather than generating complete papers autonomously.

Scholarly writing evaluation studies~\cite{yuan2024llm} show LLMs can assess paper quality. Our system goes further: not just evaluating papers but generating them from development artifacts.

\subsection{Meta-Research and Replication Studies}

Meta-research examines the research process itself~\cite{ioannidis2005most}. Replication studies~\cite{juristo2010replication} validate prior findings by repeating experiments. Our approach automates meta-research: analyzing development history to generate papers about the development process.

Software engineering has a tradition of tool papers describing novel systems~\cite{ernst2014tool}. These are typically written manually by tool authors. We automate this process, generating tool papers from commit history and design documents.

\subsection{Self-Documenting Systems}

Knuth's literate programming~\cite{knuth1984literate} interleaves code and documentation in a single source. The system generates both executable code and readable documentation from the same input.

Our approach inverts this: instead of embedding documentation in code, we extract documentation from code history. This enables retroactive documentation: generating papers about systems after they're built, using commits and development logs as evidence.

\subsection{Program Synthesis and Code Generation}

Program synthesis research~\cite{gulwani2017program} generates code from specifications. Recent LLM-based systems~\cite{chen2021codex,austin2021program} generate code from natural language descriptions.

Our system performs the inverse transformation: generating natural language (research papers) from code and development artifacts. This is conceptually similar to code summarization~\cite{haiduc2010supporting} but at paper scale rather than function scale.

\subsection{Gaps in Existing Work}

No prior work combines topic discovery from development artifacts, evidence extraction across multiple sources (commits, waves, docs), and complete research paper generation with substantive content. Existing systems either:
\begin{itemize}
\item Generate code-level documentation (Javadoc, Copilot) without paper-level narrative
\item Extract data for manual analysis (MSR) without generating papers
\item Assist human writing (Elicit, TLDR) without autonomous generation
\end{itemize}

Our contribution is end-to-end automation from artifacts to papers, with self-referential validation demonstrating the system can document itself.